% part: intuitionistic-logic
% chapter: propositions-as-types
% section: normalization

\documentclass[../../../include/open-logic-section]{subfiles}

\begin{document}

\olfileid[hi]{int}{pty}{nor}

\olsection{प्रसामान्यीकरण}

इस अनुभाग में हम सिद्ध करते हैं कि किसी अपचयन क्रम के अनुसार प्रत्येक निगमन
को एक प्रसामान्य निगमन में अपचयित किया जा सकता है; इसे
\emph{प्रसामान्यीकरण गुणधर्म} कहते हैं। हम प्रकारों के रूप में प्रतिज्ञप्तियों
की अनुरूपता का उपयोग करेंगे: हम दिखाते हैं कि प्रत्येक प्रमाण पद को एक
प्रसामान्य रूप में अपचयित किया जा सकता है; इससे प्राकृतिक निगमन की
!!{derivation}s के लिए प्रसामान्यीकरण प्राप्त होता है।

पहले हम पदों की जटिलता मापने वाले कुछ फलन परिभाषित करते हैं।
किसी !!a{formula}s की \emph{लंबाई}~$\len{!A}$ इस प्रकार परिभाषित है:
\begin{align*}
  \len{p} &= 0\\
  \len{!A \land !B} &= \len{!A} + \len{!B} + 1\\
  \len{!A \lor !B} &= \len{!A} + \len{!B} + 1 \\
  \len{!A \lif !B} &= \len{!A} + \len{!B} + 1.
\end{align*}
किसी अपचय्य~$M$ की जटिलता उसकी \emph{कट कोटि}~$\cutrank{M}$ से मापी जाती है:
\begin{align*}
  \cutrank{(\lambd[\typeof{x}{!A}][\typeof{N}{!B}])Q} &=
  \len{!A} + \len{!B} + 1 \\
  \cutrank{\ande{i}{\andi{\typeof{M}{!A}}{\typeof{N}{!B}}}} &=
  \len{!A} + \len{!B} + 1 \\
  \cutrank{\ore{\ori{i}{!A_{i}}{\typeof{M}{!A_i}}}
    {\typeof{x_1}{!A_1}}{\typeof{N_1}{!C}}
    {\typeof{x_2}{!A_2}}{\typeof{N_2}{!C}}} &=
  \len{!A} + \len{!B} + 1
\end{align*}
किसी प्रमाण पद की जटिलता उसमें उपस्थित सबसे जटिल अपचय्य से मापी जाती है,
और यदि वह प्रसामान्य है तो उसका मान $0$ होता है:
\[
  \maxrank{M} = \max \{\cutrank{N} | N \text{ पद } M \text{ का उपपद और अपचय्य है}\}
\]

\begin{lem}\ollabel{lem:subst}
  यदि $\Subst{M}{\typeof{N}{!A}}{\typeof{x}{!A}}$ एक अपचय्य है और $M
  \not\ident x$, तो निम्नलिखित में से कोई एक स्थिति होती है:
  \begin{enumerate}
  \item $M$ स्वयं एक अपचय्य है, अथवा
  \item $M$ का रूप $\ande{i}{x}$ है और $N$ का रूप
    $\andi{P_1}{P_2}$ है,
  \item $M$ का रूप $\ore{i}{x_1}{P_1}{x_2}{P_2}$ है और $N$ का रूप
    $\ori{i}{}{Q}$ है,
  \item $M$ का रूप $x Q$ है और $N$ का रूप
    $\lambd[x][P]$ है।
  \end{enumerate}
  पहली स्थिति में $\cutrank{\Subst{M}{N}{x}} = \cutrank{M}$; अन्य
  स्थितियों में $\cutrank{\Subst{M}{N}{x}} = \len{!A})$।
\end{lem}
\begin{proof}
  $M$ पर आगमन द्वारा प्रमाण।
  \begin{enumerate}
  \item यदि $M$ एक अकेला चर $y$ है और $y \not\ident x$, तो
    $\Subst{y}{N}{x}$ का मान $y$ है, अतः वह अपचय्य नहीं है।
  \item यदि $M$ का रूप $\andi{N_1}{N_2}$, या $\lambd[x][N]$, या
    $\ori{i}{!A}{N}$ है, तो $\Subst{M}{\typeof{N}{!A}}{\typeof{x}{!A}}$
    का भी वही रूप है और इसलिए वह अपचय्य नहीं है।
  \item यदि $M$ का रूप $\ande{i}{P}$ है, तो हम दो स्थितियाँ देखते हैं।
    \begin{enumerate}
      \item यदि $P$ का रूप $\andi{P_1}{P_2}$ है, तो $M \ident
        \ande{i}{\andi{P_1}{P_2}}$ एक अपचय्य है और स्पष्टतः
        \[
        \Subst{M}{N}{x} \ident
        \ande{i}{\andi{\Subst{P_1}{N}{x}}{\Subst{P_2}{N}{x}}}
        \]
        भी एक अपचय्य है। दोनों की कट कोटियाँ समान हैं।
      \item यदि $P$ एक अकेला चर है, तो प्रतिस्थापन को अपचय्य बनाने के लिए
        वह $x$ ही होना चाहिए, और $N$ का रूप $\andi{P_1}{P_2}$ होना चाहिए।
        अब
        \[
        \Subst{M}{N}{x} \ident \Subst{\ande{i}{x}}{\andi{P_1}{P_2}}{x}
        \]
        पर विचार करें, जो $\ande{i}{\andi{P_1}{P_2}}$ है। इसकी कट कोटि
        $\cutrank{x}$ के बराबर है, जो $\len{!A}$ है।
    \end{enumerate}
  \end{enumerate}
  $\ore{N}{x_1}{N_1}{x_2}{N_2}$ और $P Q$ की स्थितियाँ इसी प्रकार हैं।
\end{proof}

\begin{lem}
  यदि $M$ संकुचित होकर $M'$ बनता है, और $M$ के प्रत्येक उचित अपचय्य उपपद
  $N$ के लिए $\cutrank{M} > \cutrank{N}$ है, तो $\cutrank{M} >
  \maxrank{M'}$।
\end{lem}

\begin{proof}
  स्थितियों के अनुसार प्रमाण।
  \begin{enumerate}
  \item यदि $M$ का रूप $\ande{i}{\andi{M_1}{M_2}}$ है, तो $M'$ का मान
    $M_i$ है; चूँकि $M_i$ का प्रत्येक उपपद $M$ का भी उचित उपपद है,
    इसलिए कथन सिद्ध होता है।
  \item यदि $M$ का रूप
    $(\lambd[\typeof{x}{!A}][N])\typeof{Q}{!A}$ है, तो $M'$ का मान
    $\Subst{N}{\typeof{Q}{!A}}{\typeof{x}{!A}}$ है। $M'$ में किसी अपचय्य
    पर विचार करें। या तो $N$ में समान कट कोटि वाला अनुरूपी अपचय्य है,
    जिसकी कट कोटि मान्यता के अनुसार $\cutrank{M}$ से कम है, या उसकी कट कोटि
    $\len{!A}$ के बराबर है, जो परिभाषा से
    $\cutrank{(\lambd[x^{!A}][N])Q}$ से कम है।
  \item यदि $M$ का रूप
    \[
    \ore{\ori{i}{}{\typeof{N}{!A_i}}}
        {\typeof{x_1}{!A_1}}{\typeof{N_1}{!C}}
        {\typeof{x_2}{!A_2}}{\typeof{N_2}{!C}}
    \]
    है, तो $M' \ident \Subst{N_i}{N}{\typeof{x_i}{!A_i}}$। $M'$ में किसी
    अपचय्य पर विचार करें। या तो $N_i$ में समान कट कोटि वाला अनुरूपी अपचय्य
    है, जिसकी कट कोटि मान्यता के अनुसार $\cutrank{M}$ से कम है; या उसकी कट
    कोटि $\len{!A_i}$ के बराबर है, जो परिभाषा से
    $\cutrank{\ore{\ori{i}{}{\typeof{N}{!A_i}}}
      {\typeof{x_1}{!A_1}}{\typeof{N_1}{!C}}
      {\typeof{x_2}{!A_2}}{\typeof{N_2}{!C}}}$ से कम है।
  \end{enumerate}
\end{proof}

\begin{thm}
  सभी प्रमाण पद अपचयित होकर प्रसामान्य रूप में आते हैं; सभी !!{derivation}s
  अपचयित होकर प्रसामान्य !!{derivation}s में आती हैं।
\end{thm}

\begin{proof}
  दूसरा कथन पहले से प्राप्त होता है। हम पहले कथन को प्रमाण पद $M$ के लिए
  $m=\maxrank{M}$ पर पूर्ण आगमन द्वारा सिद्ध करते हैं।
  \begin{enumerate}

  \item यदि $m=0$, तो $M$ पहले से प्रसामान्य है।
  \item अन्यथा हम $n$ पर आगमन करते हैं, जहाँ $n$, $M$ में कट कोटि $m$ वाले
    अपचय्यों की संख्या है।
    \begin{enumerate}
    \item यदि $n=1$, तो ऐसा कोई अपचय्य $N$ चुनें कि प्रत्येक उचित उपपद $P$
      के लिए, जो स्वयं भी अपचय्य है, $m = \cutrank{N} > \cutrank{P}$ हो।
      ऐसा अपचय्य अवश्य है, क्योंकि प्रत्येक पद के उपपदों की संख्या सीमित है।

      $N$ के अपचयफल को $N'$ लिखें। अब प्रमेयिका से
      $\maxrank{N'} < \maxrank{N}$; अतः कट कोटि $m$ वाले अपचय्यों की संख्या
      $n$ घट जाती है। इसलिए $m$ घटता है ($1$ या अधिक से), और हम $m$ के लिए
      आगमन परिकल्पना लागू कर सकते हैं।
    \item आगमन चरण के लिए मान लें कि $n>1$। प्रक्रिया समान है, केवल $n$
      घटकर कोई धनात्मक संख्या रह जाता है और इसलिए $m$ नहीं बदलता। हम बस
      $n$ के लिए आगमन परिकल्पना लागू करते हैं।
    \end{enumerate}
\end{enumerate}
\end{proof}

प्रकारित $\lambd$-पदों के प्रसामान्यीकृत होने का अर्थ है कि प्रत्येक ऐसे
पद का \emph{एक प्रसामान्य रूप होता है}। यदि हम $\lambd$-पदों को उनके
प्रसामान्य रूप तक $\beta$-अपचयित करने को किसी गणना का निष्पादन और प्रसामान्य
रूप को उस गणना का मान समझें, तो इसका अर्थ है कि प्रकारित $\lambd$-कलन की
प्रत्येक गणना अंततः रुकती है और कोई मान उत्पन्न करती है। इसके अतिरिक्त,
प्रकार-संरक्षण सुनिश्चित करता है कि उस मान का प्रकार सही है। उदाहरणतः, यदि
$N$ का प्रकार $!A \lif !B$ है (अर्थात वह ऐसा फलन है जिसे प्रकार~$!A$ के
निवेशों पर काम करके प्रकार~$!B$ के मान उत्पन्न करने चाहिए) और हम उसे
प्रकार~$!A$ के पद~$M$ (निवेश) पर लागू करें, तो $(NM)$ अपचयित होकर पद $N'$
(निर्गम) बनता है और यह प्रत्याभूत है कि उस निर्गम का प्रकार~$!B$ है।

वास्तव में, प्रकारित $\lambd$-कलन के पद केवल उपपदों पर $\beta$-अपचयन लागू
करने के किसी एक ढंग से ही प्रसामान्य रूपों तक नहीं पहुँचते, बल्कि
$\beta$-अपचयन लागू करने का \emph{हर} ढंग अंततः किसी प्रसामान्य रूप पर समाप्त
होता है। इस गुणधर्म को \emph{प्रबल प्रसामान्यीकरण} कहते हैं। इसका अर्थ है कि
न केवल गणना के अंततः रुकने और कोई मान देने की प्रत्याभूति का कोई ढंग है,
बल्कि गणना को निष्पादित करने का \emph{हर} ढंग अंततः कोई मान देता है।

हम कैसे जानें कि गणना को निष्पादित करने के विभिन्न ढंग \emph{एक ही} मान देते
हैं? अभी तक हम (प्रकार-संरक्षण से) केवल यह जानते हैं कि उससे मिलने वाले सभी
मानों का प्रकार समान है। प्रकारित $\lambd$-कलन का एक और महत्त्वपूर्ण गुणधर्म
है: वह \emph{संगमशील}, अथवा \emph{चर्च--रॉसर}, है। यदि $N \red N_1$ और
$N \red N_2$, तो कोई पद~$N'$ ऐसा है कि $N_1 \red N'$ और $N_2 \red N'$।
याद करें कि $N \red N'$ का अर्थ है: $N'$ शून्य या अधिक $\redone$-अपचयन
चरणों से प्राप्त होता है। यदि $N_1$ प्रसामान्य रूप में है, तो $N_1 \red N'$
को संतुष्ट करने वाला एकमात्र पद $N'$ स्वयं $N_1$ है (शून्य $\redone$-अपचयन
चरणों द्वारा)। अतः यदि $N_1$ और $N_2$ दोनों प्रसामान्य रूप में हैं, तो
$N_1 = N' = N_2$। दूसरे शब्दों में, क्योंकि $\red$ प्रबल रूप से
प्रसामान्यीकृत होता है, किसी पद को अपचयित करने का हर ढंग अंततः प्रसामान्य
रूप देता है। क्योंकि $\red$ चर्च--रॉसर है, कोई भी दो प्रसामान्य रूप समान
हैं। इसलिए प्रकारित $\lambd$-कलन में गणना सदैव समाप्त होती है और गणना चाहे
जिस ढंग से आगे बढ़े, उसका मान (प्रसामान्य रूप) समान रहता है।

किसी संबंध को संगमशील सिद्ध करना सामान्यतः कठिन होता है। फिर भी हम निम्न
दुर्बलतर प्रतिबंध सरलता से स्थापित कर सकते हैं:

\begin{prop}[दुर्बल चर्च--रॉसर गुणधर्म]
यदि $N \redone N_1$ और $N \redone N_2$, तो कोई पद $N'$ ऐसा है कि
$N_1 \red N'$ और $N_2 \red N'$।
\end{prop}

\begin{proof}
  दोनों पद $N_1$ और $N_2$, $N$ में क्रमशः अपचय्य $M_1$ अथवा $M_2$ को उसके
  अपचयफल से बदलने पर मिलते हैं। अपचय्य $M_1$ और $M_2$, $N$ के उपपद हैं और
  इसलिए परस्पर आच्छादित नहीं हो सकते। इसका अर्थ है कि या तो एक दूसरे का
  उपपद है, या वे $N$ में अलग-अलग, अनाच्छादित स्थानों पर आते हैं। बाद वाली
  स्थिति देखें: $N$ का रूप $N(M_1,M_2)$ है। $M_1$ को उसके अपचयफल $M_1'$ से
  बदलने पर $N_1$, अर्थात $N(M_1',M_2)$, मिलता है; और $M_2$ को उसके अपचयफल
  $M_2'$ से बदलने पर $N_2$, अर्थात $N(M_1,M_2')$, मिलता है। दोनों अपचयित
  होकर $N(M_1',M_2')$ बनते हैं। दूसरी स्थिति में मान लें कि $M_2$, $M_1$ का
  उपपद है (दूसरी दिशा सममित है), और $M_1$ किस प्रकार का अपचय्य है इसके
  अनुसार स्थितियाँ अलग करें। उदाहरणतः यदि $M_1 \ident
  \proj{1}{\pair{O_1}{O_2}}$, तो वह अपचयित होकर $O_1$ बनता है। यदि $M_2$,
  $O_1$ का उपपद है, तो $M_2$ के रूपांतरण से $O_1'$ मिलता है। अतः पहले $M_1$
  और फिर $M_2$ को रूपांतरित करने पर $O_1'$ मिलता है। किंतु पहले $M_2$ को
  रूपांतरित करने पर $\proj{1}{\pair{O_1'}{O_2}}$ मिलता है और यह भी अपचयित
  होकर~$O_2$ बनता है।
\end{proof}

\begin{lem}[न्यूमन की प्रमेयिका]
यदि $\red$ प्रबल रूप से प्रसामान्यीकृत होता है और उसमें दुर्बल चर्च--रॉसर
गुणधर्म है, तो वह संगमशील है।
\end{lem}

\begin{proof}
  क्योंकि $\red$ प्रबल रूप से प्रसामान्यीकृत होता है, प्रत्येक अपचयन अनुक्रम
  किसी प्रसामान्य रूप में समाप्त होता है। प्रसामान्य रूप अद्वितीय हैं यदि और
  केवल यदि $\red$ संगमशील है। अतः मान लें कि किसी पद $N$ के दो भिन्न
  प्रसामान्य रूप $N'$ और $N''$ हैं, जो इन अपचयन अनुक्रमों से मिलते हैं:
  \begin{align*}
  N & = N_1' \redone N_2' \redone \dots \redone N_k' = N' \text{ और}\\
  N & = N_1'' \redone N_2'' \redone \dots \redone N_l'' = N''
  \end{align*}
  (अपचयन अनुक्रम की लंबाई $>0$ होनी चाहिए, क्योंकि अन्यथा $N$ पहले से
  प्रसामान्य रूप में होता।)

  यदि $N_1' = N_1''$, तो उसके दो भिन्न प्रसामान्य रूप ($N'$ और $N''$) हैं।
  यदि $N_1' \neq N_1''$, तो दुर्बल चर्च--रॉसर गुणधर्म से कोई $N'''$ ऐसा है
  कि $N_1' \red N'''$ और $N_1'' \red N'''$। क्योंकि $N' \neq N''$, या तो
  $N''' \neq N'$ अथवा $N''' \neq N''$। फलतः $N_1'$ या $N_1''$ में से किसी
  एक के दो भिन्न प्रसामान्य रूप हैं। हमने दिखाया है कि यदि $N$ के दो भिन्न
  प्रसामान्य रूप हैं, तो कोई $N_1$ ऐसा है कि $N \redone N_1$ और $N_1$ के
  दो भिन्न प्रसामान्य रूप हैं। स्पष्टतः हम इसी प्रकार आगे बढ़कर
  $N \redone N_1 \redone N_2 \redone \dots$ बना सकते हैं, किंतु यह असंभव
  है क्योंकि $\redone$ प्रबल रूप से प्रसामान्यीकृत होता है।
\end{proof}


\end{document}

